\documentclass[runningheads]{llncs}

\usepackage[T1]{fontenc}
\usepackage{amsmath,amssymb}
\usepackage{graphicx}
\usepackage{booktabs}
\usepackage{url}
\usepackage[hidelinks]{hyperref}
\renewcommand\UrlFont{\rmfamily}
\setlength{\textfloatsep}{10pt plus 2pt minus 2pt}
\renewcommand{\topfraction}{0.92}
\renewcommand{\bottomfraction}{0.75}
\renewcommand{\textfraction}{0.06}
\renewcommand{\floatpagefraction}{0.85}
\setlength{\intextsep}{8pt plus 2pt minus 2pt}

\begin{document}

\title{The Accuracy Illusion: How Overlapping Samples Manufacture
Confidence, Not Skill, in Intraday Index Forecasting}

\titlerunning{The Accuracy Illusion}

\author{Rachit Ranka\inst{1} \and Aaditya Saini\inst{1}}

\authorrunning{R. Ranka and A. Saini}

\institute{Department of Computer Science Engineering, Chandigarh University,
Mohali, India\\
\email{rankarachit5@gmail.com}, \email{aadityasaini899@gmail.com}}

\maketitle

\begin{abstract}
Intraday forecasting systems routinely report directional hit-rates between
60\% and 85\%, and those figures move real capital. We show they are not valid
estimates of skill. A $k$-bar-ahead forecast issued every bar shares $k{-}1$ of
its $k$ outcome bars with its neighbour; when the model's calls also persist, as
they do on a live system we instrument, the reported hit-rate acquires a
variance far larger than its nominal sample size implies. The reading is not
biased upward. Its stated precision is. A forecaster with \emph{exactly zero
skill}, simulated against real NSE index bars, posts a $\geq$71\% session once
in 48 while reporting an interval $1.68\times$ too narrow, rising to
$3.40\times$ at the per-minute density these systems publish. We give the closed form behind this, the classical design effect with clusters
given by runs of identical calls, and a calibrated interval built on it that
covers at 93\% where the conventional one covers at 73\%. Applied to our own
results it leaves nothing: across 48 model--index--horizon tests, 16 clear
$p<0.05$ scored per bar and \emph{zero} clear it per event. Break-even accuracy
derived from measured move distributions is 71.5\% at thirty minutes against a
best achieved 51.1\%, and at five minutes no accuracy suffices under a symmetric
bracket. Code and raw results are public.

\keywords{Overlapping observations \and Backtest overfitting \and Directional
forecasting \and Negative results \and Reproducibility \and Index derivatives}
\end{abstract}

\section{Introduction}

A system reporting 71\% directional accuracy is, on its face, extraordinary.
Work that controls for scoring unit, costs and multiple testing places realistic
out-of-sample intraday index accuracy close to a coinflip
\cite{timmermann2004,sullivan1999}, while surveys of the applied literature
catalogue many higher figures reported without those controls \cite{sezer2020}.

This paper began as an attempt to build such a system. Over fourteen months we
ran a live multi-horizon forecasting platform for Indian index derivatives, and
it reported 71\%. That number was wrong. The session producing it cannot be
re-audited, the portion of the log holding it having been lost before this
analysis; every figure below comes from sessions that survive.

What the overlapping-observations literature does not document is the magnitude
of the distortion in a now-ubiquitous practice: the every-minute multi-horizon
dashboard, where overlap meets a second ingredient, prediction persistence, and
is amplified across a panel.

\paragraph{Contributions.}
(i) We measure the artefact as it occurs, showing that overlap alone changes
nothing and that inflation requires overlap \emph{and} persistent predictions,
with persistence calibrated from a production log rather than assumed. (ii) We
derive the effective sample size of such a stream, show it is governed by
$\mathbb{E}[L^2]/\mathbb{E}[L]$ over run lengths, and calibrate an interval that
covers at its nominal rate with no fitted constant. (iii) We apply it to our own
results and to a complete control set --- shuffled labels, injected positive
controls, eight architectures, a block-bootstrap null --- and report that nothing
survives. (iv) We invert accuracy into a required-versus-achieved framing that
shows the gap is not close, and is in places unbridgeable at any accuracy. Code,
raw results and the free-data corpus are at
\url{https://github.com/rachitranka25/accuracy-illusion}.

\section{Related Work}

Hansen and Hodrick \cite{hansen1980} characterise inference under overlapping
returns and Newey and West \cite{newey1987} give the standard correction, but
that work concerns
regression coefficients. We apply the same logic to what practitioners publish,
the session hit-rate, and reach the classical design effect
\cite{kish1965} with clusters given by runs of identical calls. The
contribution is the application to forecast streams, not the statistic.
Boudoukh et al.\ \cite{boudoukh2008} are the closest antecedent: they show
overlap manufacturing apparent long-horizon predictability that vanishes under
correct inference.

On multiplicity, Harvey et al.\ \cite{harvey2016} argue the conventional $t>2$
threshold is indefensible after decades of factor mining. Two deflation
statistics address the concern from different directions: one discounts a
reported Sharpe ratio by the number of trials behind it \cite{bailey2014}, the
other estimates how often an in-sample winner fails out of sample
\cite{bailey2017}. White's reality check \cite{white2000} and the
false-discovery-rate procedure of Benjamini and Hochberg \cite{bh1995} are
complementary, and Sullivan et al.\ \cite{sullivan1999} apply the former to a
universe of technical trading rules --- the direct precedent for our strategy
sweep. On machine learning for finance, Fischer and Krauss \cite{fischer2018}
report performance gross of costs and Gu et al.\ \cite{gu2020} modest
out-of-sample $R^2$, while Zhang et al.\ \cite{zhang2019} and Kolm et al.\
\cite{kolm2023} find genuine short-horizon alpha from message-level data rather
than bars.

The point is not that any of that work is wrong but that the scoring unit
differs: those studies aggregate per day or per month over non-overlapping
periods, where the artefact cannot arise, or score per event. Per-bar scoring of
a multi-bar horizon, the condition under which it does arise, is characteristic
of intraday dashboards rather than of the published work we cite. None of them
reports a multiple-testing correction or a shuffled-label null in the form used
here. This paper scores per event, corrects by Bonferroni and
Benjamini--Hochberg, carries a shuffled-label null and injected positive
controls, models costs explicitly, and releases code and raw output.

\section{Data and Method}

All data is obtainable free from a retail brokerage account and public exchange
archives: 26\,265 and 26\,267 five-minute bars on BANKNIFTY and NIFTY~50 from
February 2025 to July 2026, 9.6k--14k bars on three further indices, 1\,382\,738
rows of end-of-day option chain with open interest, a ten-year chain history
covering 2\,181 days, 515 days of participant-wise open interest and 639 days of
index futures carrying real volume. One practical note: index spot prints zero
volume, so a feature computed as a ``VWAP'' on spot is an unweighted expanding
mean, and real volume exists only on futures.

Every accuracy below passes through one harness. Splits are expanding-window
walk-forward with the $k$ bars whose labels overlap the test window purged from
training \cite{prado2018}. Accuracies carry Clopper--Pearson \cite{clopper1934}
intervals, exact binomial $p$-values and the Pesaran--Timmermann statistic
\cite{pt1992}. Every model is additionally trained on randomly permuted labels;
if real and shuffled scores coincide, the pipeline is reading noise. A synthetic
feature of known strength is injected as a positive control, since a pipeline
that cannot recover planted signal cannot be trusted to report a null. Strategy
returns use non-overlapping $k$-bar holds with explicit costs.

\paragraph{Units.}
All analysis is on 5-minute bars, so the 30-minute horizon is $k=6$ and a
session yields $n=69$ forecasts. The production system emits a forecast every
\emph{minute}; pooled, its run lengths average 15.2 minutes with a longest of
263, and the pooled mean run in 5-minute bars is 3.4.

\section{The Mechanism}
\label{sec:mech}

\subsection{Overlap is arithmetic}

Measure the autocorrelation of the $k$-bar forward return two ways. On
overlapping windows it reaches 0.841 at lag 1 and decays to approximately zero at
\emph{exactly} lag $k$; sampled so that no two observations share a bar it is
0.067 at lag 1 and indistinguishable from zero thereafter. The pattern reproduces
at $k=6$ and $k=12$, with the decay landing on lag 6 and lag 12 respectively. The
correlation is a property of the sampling scheme, not of the market, and it is
the whole of the mechanism below.

\subsection{Persistence is the second ingredient}

Overlap alone is insufficient, and getting this wrong was instructive: if
predictions are drawn independently each bar the hit indicator is i.i.d.\
Bernoulli however autocorrelated the outcome, and our first simulation made that
error and correctly found nothing. Real models do not re-decide every bar.
Measured on the production log at the 30-minute horizon, BANKNIFTY calls persist
for a mean of 25.9 minutes and NIFTY~50 for 10.3, against a mean run of one for a
model that re-drew each bar. When both call and outcome persist, one lucky
trending stretch is recorded as dozens of hits.

\subsection{What a zero-skill forecaster reports}

We simulate a fair coin held for a run length drawn from the measured
distribution, evaluated against real bars and scored as a dashboard scores it
(Table~\ref{tab:zero}).

\begin{table}[htbp]
\caption{A forecaster with exactly zero skill, 30-minute horizon, scored per bar
versus per event. BANKNIFTY; NIFTY 50 agrees within 0.6 points on every row.}
\label{tab:zero}
\centering
\footnotesize
\begin{tabular}{@{}lrr@{}}
\toprule
 & \textbf{Per bar} & \textbf{Per $k$ bars} \\
\midrule
Reported sample size per session & 69 & 12 \\
Stated binomial 95\% CI & $\pm$11.8 pts & $\pm$28.3 pts \\
Realised spread across sessions & $\pm$19.9 pts & $\pm$28.0 pts \\
\textbf{Understatement factor} & \textbf{1.68$\times$} & \textbf{0.99$\times$} \\
\midrule
$P(\text{session} \geq 60\%)$ & 16.2\% & 18.9\% \\
$P(\text{session} \geq 71\%)$ & 2.1\% & 7.1\% \\
\bottomrule
\end{tabular}
\end{table}

The table is easily read backwards. Per-bar scoring does \emph{not} produce
extreme readings more often than honest scoring; it produces them \emph{less}
often, and reports them with an interval $1.68\times$ too narrow because the
quoted $n$ of 69 is a fiction. Overlap manufactures unwarranted
\emph{confidence}, not high hit-rates.

\subsection{At the real granularity, and across a panel}

The production system forecasts every minute, so a session yields about 313
forecasts rather than 69, and simulating the same zero-skill forecaster against
the log's own outcome series raises the understatement factor to $3.40\times$
while extreme sessions become slightly \emph{rarer}. Every 5-minute-bar figure
here therefore understates the artefact. A dashboard also watches more than one
number: two indices at four horizons is eight hit-rates per session, and
simulating the panel jointly gives $P(\exists\,\text{series}\geq 71\%) = 17.1\%$
--- about one session in six, so a headline 71\% arrives within the first weeks
of live operation with no edge at all.

\section{A Calibrated Estimator}
\label{sec:est}

\subsection{Effective sample size}

Let $h_t = \mathbf{1}\{\hat{y}_t = y_t\}$ and $\hat{p} = n^{-1}\sum_t h_t$.
Practitioners quote $\mathrm{Var}(\hat{p}) = p(1-p)/n$, valid only for independent
$h_t$. In general
\begin{equation}
\mathrm{Var}(\hat{p}) = \frac{p(1-p)}{n}
\Big[1 + 2\sum_{j=1}^{n-1}\big(1 - \tfrac{j}{n}\big)\rho_j\Big],
\label{eq:var}
\end{equation}
with $\rho_j = \mathrm{Corr}(h_t,h_{t+j})$. Because the dependence is
block-shaped rather than smoothly decaying, this collapses: if the sequence
decomposes into runs of length $L_i$ within which $h_t$ is constant, then
\begin{equation}
n_{\text{eff}} = \frac{n}{\sum_i L_i^2/n}
\;\xrightarrow[n\to\infty]{}\; \frac{n\,\mathbb{E}[L]}{\mathbb{E}[L^2]}.
\label{eq:neff}
\end{equation}
This is the classical design effect \cite{kish1965}; what is new is that a
forecast stream has exactly this structure and that nobody reports it, and that
the inflation is governed by $\mathbb{E}[L^2]$, so the tail matters.

\subsection{Validation by coverage}

Table~\ref{tab:cov} reports coverage against a forecaster of known skill whose
calls persist per the measured distribution, with true skill measured from the
generative process rather than assumed.

\begin{table}[htbp]
\caption{Coverage of a nominal 95\% interval on a session hit-rate. BANKNIFTY,
30-minute horizon, nominal $n=69$; NIFTY 50 agrees within 2 points. HAC-L gives
the Bartlett kernel a run-structure bandwidth; \emph{derived} divides the design
effect by the alternation factor of \eqref{eq:alt}.}
\label{tab:cov}
\centering
\footnotesize
\begin{tabular}{@{}lcccccr@{}}
\toprule
\textbf{True skill} & \textbf{Naive} & \textbf{HAC} & \textbf{HAC-L} &
\textbf{Design} & \textbf{Derived} & $\boldsymbol{n_{\text{eff}}}$ \\
\midrule
50\% & 73.4 & 86.2 & 88.0 & 99.7 & \textbf{93.2} & 11.9 \\
54\% & 73.9 & 86.6 & 88.1 & 99.6 & \textbf{93.5} & 11.7 \\
59\% & 73.0 & 85.2 & 86.8 & 99.6 & \textbf{93.5} & 10.9 \\
\bottomrule
\end{tabular}
\end{table}

HAC and a block bootstrap (86.2\%, omitted from the table) both under-cover by
about nine points. The diagnosis is bandwidth: the automatic Bartlett rule gives
3.7 bars at $n=69$ against a pooled mean run of 3.4 and a tail reaching 53, and a
run-structure bandwidth recovers only one to two points. The design effect \eqref{eq:neff} removes the failure mode entirely and over-corrects in the
other direction, reaching 99.7\%. That over-statement is derivable rather than
incidental: treating runs as independent blocks ignores that a binary sequence's
runs alternate deterministically, a run of hits always being followed by a run of
misses. For geometric run lengths of mean $\mu$,
\begin{equation}
\frac{\mathbb{E}[L^2]}{\mathbb{E}[L]} = 2\mu - 1,
\qquad
\text{whereas exactly } \frac{1+\rho}{1-\rho} = \mu - 1,
\label{eq:alt}
\end{equation}
so \eqref{eq:neff} over-states the dependence by a factor $(2\mu-1)/(\mu-1)$, a
quantity that tends to 2 for long runs. Evaluated per session it averages 2.44
at a mean hit-run length of $\mu = 3.53$; the residual is the departure of real
run lengths from geometry.

Dividing by that factor, computed from each sample's own mean run length,
introduces no free constant and gives 93.1--93.5\% coverage across both indices
and all three skill levels. We report it as the recommended estimator, and the
raw design effect as the conservative bound.

\subsection{Reliability, and a real session}

Equation~\eqref{eq:neff} depends on a tail-driven second moment estimated from a
few dozen runs, and a tail is what a short sample estimates worst: against a
pooled reference of 10.6, a single session's $n_{\text{eff}}$ lands between
$0.63\times$ and $1.73\times$ that value nine times in ten, so it should be
estimated from a run of sessions. That matters immediately. On the production log
the best NIFTY~50 session, 69.6\% over 335 forecasts, has a naive interval
excluding 50\% and corrects to $[50.0, 89.1]$, which still excludes it --- but at
the 5th percentile of the $n_{\text{eff}}$ uncertainty it becomes
$[45.1, 94.0]$ and no longer does. The margin is smaller than the estimator's own
noise, so the session supports no claim of skill. An earlier draft reported that
it did.

Per-event scoring is unbiased and is what we use for every backtest here, but it
is insufficient for a live stream: a dashboard cannot wait $k$ bars, subsampling
throws away 83\% of the observations, and which of the $k$ offsets is chosen
moves the same session by 7 to 20 points.

\section{Predictability Under Correct Scoring}
\label{sec:pred}

\subsection{Is there anything to predict?}

Before asking which model is best, we ask whether any model could succeed. An
unconstrained forest was fitted until it memorised the training data, on five
indices across two exchanges. Training accuracy reaches 100.0\% everywhere and
test accuracy stays at chance: 50.5 and 50.0 on BANKNIFTY and NIFTY~50 over
16\,856 and 16\,853 scored bars, 49.7--50.2 on the other three. Permuting labels
moves accuracy by 0.3 to 0.7 points, no real-label result is distinguishable from
a coinflip ($p = 0.22$ to $0.91$), and both injected positive controls, planted
at about 55\% and 65\%, are recovered at $p<0.001$ everywhere. The null is a
property of the data, not of the estimator.

\subsection{Eight model families, three horizons}

Eight families were trained on identical splits, features, seeds and controls:
logistic regression, gradient boosting, LightGBM \cite{ke2017}, XGBoost
\cite{chen2016}, $k$-nearest neighbours, an LSTM \cite{hochreiter1997}, a
Transformer \cite{vaswani2017} and a compact Temporal Fusion Transformer
\cite{lim2021}. Counting families significant at $p<0.05$, per-bar scoring yields
0 and 2 at 15 minutes on BANKNIFTY and NIFTY~50, 3 and 1 at 30 minutes and 5 and
5 at 60: 16 of 48 tests. At an effective sample size, and again scored one
observation per event, the count is \emph{zero} in all six cells.

Scored per bar, XGBoost reaches 51.3\% on BANKNIFTY at $p=0.001$ --- in an
earlier draft our one surviving positive result, described as robust to
Bonferroni correction. It does not survive our own correction: re-scored one
observation per 6 bars it is 50.4\% at $p=0.692$, and widening intervals to
$n_{\text{eff}}$ puts 50\% inside every one. The count of \emph{apparent}
discoveries rises with horizon, exactly as it should if the artefact scales with
$k$; the count of surviving discoveries stays at zero.

The zoo contains its own diagnostic: the NIFTY~50 Transformer sits at 48.4\%
with $p<0.001$, which at the nominal $n$ is 4.3 standard errors \emph{below}
chance, where the largest $|z|$ across 48 tests should land near 2.9--3.2. That
is the variance inflation of Section~\ref{sec:mech} surfacing in our own results;
at $n_{\text{eff}}$ it vanishes.

\subsection{Eighty-six strategies}

Separately, 86 parameterised rules --- momentum, mean reversion, moving-average
crosses, RSI, Bollinger, breakout, gap, volatility-regime, time-of-day and streak
variants --- were evaluated on non-overlapping 30-minute holds at 4\,bp cost with
a one-sided test for profit. Two have a positive Sharpe ratio; none is
significantly profitable, before or after Benjamini--Hochberg.

The deflated Sharpe ratio returns an implausible null benchmark of 7.77 here,
because trial Sharpe dispersion is dominated by deterministic turnover
differences rather than by noise, violating the assumption it rests on. We
replaced it with a circular block bootstrap \cite{politis1992} that resamples
forward returns while holding each strategy's positions fixed: the best observed
Sharpe is 0.21 against a null best-of-86 median of 0.44, with
$P(\text{null} \geq \text{observed}) = 0.658$. The winner of the sweep is worse
than what selection alone produces from no edge at all. The probability of
backtest overfitting via combinatorially symmetric cross-validation
\cite{bailey2017} is degenerate here and reported as such.

\section{Economic Significance}
\label{sec:econ}

A hit-rate without a payoff ratio is not a result. For a symmetric bracket at the
typical move $m$ with round-trip cost $c$, expectancy is $m(2p-1) - c$, so
break-even requires
\begin{equation}
p^{\ast} = \tfrac{1}{2} + \frac{c}{2m}.
\label{eq:be}
\end{equation}
The ratio $c/m$ is the entire problem. At 4\,bp round-trip cost, with the
median move taken from seventeen months of bars, $p^{\ast}$ is 79.8\%, 71.5\%
and 65.4\% at 15, 30 and 60 minutes on BANKNIFTY and 84.9\%, 75.1\% and 67.7\%
on NIFTY~50, against best per-event accuracies of 51.3\%, 51.1\% and 52.2\%, and
50.5\%, 51.2\% and 50.1\%. A hold-to-horizon rule, which collects the mean move
rather than the median, lowers the requirement to 71.3--60.8\% and 74.9--62.7\%
and does not close the gap either.

At the shortest horizon the median move is smaller than the round trip, so under
a symmetric bracket $p^{\ast}$ exceeds 100\% and no accuracy suffices. That is
specific to the bracket: moves are right-skewed, so a hold-to-horizon rule
requires 86\% and 93\% instead --- far beyond anything reported here or in the
literature we cite.

Second, instrument choice can make the requirement unreachable. On the BANKNIFTY
30-minute horizon, where the median move is 9.3\,bp, break-even is 55.4\% at
1\,bp of cost, 71.5\% at 4\,bp and 93.0\% at 8\,bp; at the 20\,bp a retail
at-the-money option round trip costs over that horizon, $c > m$ and
$p^{\ast}>1$, so a perfect forecaster still loses on the median trade. The
instrument most retail participants use, over the horizon they hold it, has no
directional accuracy that makes it profitable.

A genuine 85\% would clear break-even by 13.5 points on BANKNIFTY futures and
still lose on options: such a claim is either an artefact of the kind measured in
Section~\ref{sec:mech} or true and far less valuable than it sounds.

\section{What Is Predictable Instead}

These results went through the same harness, and what it leaves is narrow enough
to state precisely: \emph{none of this project's findings survives as a finding
of ours.}

\paragraph{Volatility, confirmed but not ours.}
A binary magnitude target reaches 70.8\% per bar and 71.3\% per event on
BANKNIFTY (72.0 and 72.1 on NIFTY~50), with shuffled controls at 52--54\%. That
it survives the correction which destroyed the directional result matters, since
an instrument flattening everything would be useless. But a persistence rule
already scores 67.4\%, the edge is $+3.3$ to $+3.7$ points per event, and at
$n_{\text{eff}}$ the two intervals overlap: this confirms volatility clustering
\cite{engle1982} rather than discovering anything.

\paragraph{Premium selling and positioning, withdrawn.}
Short at-the-money straddles held to expiry, on real end-of-day premiums across
four entry points, are not significantly profitable on either index before or
after Benjamini--Hochberg, and no bootstrap interval on the mean excludes zero on
the profitable side; recentring and resampling gives a best-of-four null median
of 17.8 against an observed 12.1 on NIFTY ($P=0.68$). The exchange also moved
BANKNIFTY options from weekly to monthly expiry mid-sample, so the indices were
never comparable: split on expiry spacing the weekly sub-sample returns $-67.5$
points per trade and the monthly $+57.5$. The premium is real and documented \cite{carr2009}; our measurement of it as an edge is withdrawn.
Option-chain positioning fares no better: over 2\,181 days, six of eight
price-residual features clear $p<0.05$ uncorrected and \emph{one} survives
Benjamini--Hochberg --- the daily change in the put--call open-interest ratio
($\text{IC}=0.059$, $t\approx 2.8$), also the only feature strengthening in the
second half. Nothing in the price-only set clears $p<0.05$ at all.

\section{Limitations}

The alternation correction assumes geometric run lengths; real runs depart from
that, which is why coverage lands near 93\% rather than 95\%. Two indices carry
the primary intraday analysis over seventeen months in one market; the ceiling
test replicates across five, but the model zoo and strategy sweep do not.
Break-even is computed for two payoff models, neither of which models the price
path. Two figures that were re-run did not reproduce and are withdrawn in place
rather than quietly deleted, and our Temporal Fusion Transformer is a compact
variant, not the published architecture \cite{lim2021}. Finally, the order-flow
hypothesis, where the literature locates genuine short-horizon alpha
\cite{kolm2023,zhang2019}, is untested here for want of message-level
history.

\section{Conclusion}

The headline metric of intraday directional forecasting is not merely noisy; its
stated precision is systematically wrong, by enough to make a coinflip look like
a discovery. Applied to our own work the correction leaves nothing standing. Sixteen of 48
model--index--horizon tests clear $p<0.05$ per bar and none clears it per event;
86 strategies produce a winner worse than a bootstrap null; premium selling,
published as this project's one measurable edge, fails the same correction; and
the one untouched result confirms volatility clustering rather than finding
anything new.

The honest summary is that this project found nothing tradeable, and its value
lies in how thoroughly that was established. What we offer is a method: an
effective sample size for overlapping persistent forecasts, a calibrated interval
built on it, and a framing that turns an accuracy figure into a question about
cost. Scoring unit, effective sample size, multiplicity correction, a
shuffled-label control and an explicit cost model are not optional refinements:
without them an accuracy figure means nothing, and we know that because we
published several that did not.

\subsubsection{Reproducibility, data availability and scope.}
Every number here is regenerated by a script in the public repository, with the
raw output of each committed alongside it, at
\url{https://github.com/rachitranka25/accuracy-illusion}. That repository holds
the project's complete research record, including the positive results withdrawn
above; all experiments use a fixed seed, and collectors for every data source are
included so the corpus can be rebuilt. This is a measurement study, not
investment advice, and its central finding is one we apply first to our own
published claims.

\scriptsize
\begin{thebibliography}{27}

\setlength{\itemsep}{0pt}
\setlength{\parskip}{0pt}
\setlength{\parsep}{0pt}

\bibitem{timmermann2004}
Timmermann, A., Granger, C.W.J.: Efficient market hypothesis and forecasting.
Int. J. Forecast. \textbf{20}(1), 15--27 (2004)

\bibitem{sullivan1999}
Sullivan, R., Timmermann, A., White, H.: Data-snooping, technical trading rule
performance, and the bootstrap. J. Finance \textbf{54}(5), 1647--1691 (1999)

\bibitem{sezer2020}
Sezer, O.B., Gudelek, M.U., Ozbayoglu, A.M.: Financial time series forecasting
with deep learning: a systematic literature review 2005--2019. Appl. Soft
Comput. \textbf{90}, 106181 (2020)

\bibitem{hansen1980}
Hansen, L.P., Hodrick, R.J.: Forward exchange rates as optimal predictors of
future spot rates: an econometric analysis. J. Polit. Econ. \textbf{88}(5),
829--853 (1980)

\bibitem{newey1987}
Newey, W.K., West, K.D.: A simple, positive semi-definite, heteroskedasticity
and autocorrelation consistent covariance matrix. Econometrica \textbf{55}(3),
703--708 (1987)

\bibitem{kish1965}
Kish, L.: Survey Sampling. Wiley, New York (1965)

\bibitem{boudoukh2008}
Boudoukh, J., Richardson, M., Whitelaw, R.F.: The myth of long-horizon
predictability. Rev. Financ. Stud. \textbf{21}(4), 1577--1605 (2008)

\bibitem{harvey2016}
Harvey, C.R., Liu, Y., Zhu, H.: \ldots and the cross-section of expected
returns. Rev. Financ. Stud. \textbf{29}(1), 5--68 (2016)

\bibitem{bailey2014}
Bailey, D.H., L\'opez de Prado, M.: The deflated Sharpe ratio: correcting for
selection bias, backtest overfitting, and non-normality. J. Portf. Manage.
\textbf{40}(5), 94--107 (2014)

\bibitem{bailey2017}
Bailey, D.H., et al.: The probability of
backtest overfitting. J. Comput. Finance \textbf{20}(4), 39--70 (2017)

\bibitem{white2000}
White, H.: A reality check for data snooping. Econometrica \textbf{68}(5),
1097--1126 (2000)

\bibitem{bh1995}
Benjamini, Y., Hochberg, Y.: Controlling the false discovery rate: a practical
and powerful approach to multiple testing. J. R. Stat. Soc. B \textbf{57}(1),
289--300 (1995)

\bibitem{fischer2018}
Fischer, T., Krauss, C.: Deep learning with long short-term memory networks for
financial market predictions. Eur. J. Oper. Res. \textbf{270}(2), 654--669
(2018)

\bibitem{gu2020}
Gu, S., Kelly, B., Xiu, D.: Empirical asset pricing via machine learning. Rev.
Financ. Stud. \textbf{33}(5), 2223--2273 (2020)

\bibitem{zhang2019}
Zhang, Z., Zohren, S., Roberts, S.: DeepLOB: deep convolutional neural networks
for limit order books. IEEE Trans. Signal Process. \textbf{67}(11), 3001--3012
(2019)

\bibitem{kolm2023}
Kolm, P.N., Turiel, J., Westray, N.: Deep order flow imbalance: extracting alpha
at multiple horizons from the limit order book. Math. Finance \textbf{33}(4),
1044--1081 (2023)

\bibitem{prado2018}
L\'opez de Prado, M.: Advances in Financial Machine Learning. Wiley, Hoboken
(2018)

\bibitem{clopper1934}
Clopper, C.J., Pearson, E.S.: The use of confidence or fiducial limits
illustrated in the case of the binomial. Biometrika \textbf{26}(4), 404--413
(1934)

\bibitem{pt1992}
Pesaran, M.H., Timmermann, A.: A simple nonparametric test of predictive
performance. J. Bus. Econ. Stat. \textbf{10}(4), 461--465 (1992)

\bibitem{ke2017}
Ke, G., et al.: LightGBM: a highly efficient gradient boosting decision tree. In: Advances in
Neural Information Processing Systems, vol. 30, pp. 3146--3154 (2017)

\bibitem{chen2016}
Chen, T., Guestrin, C.: XGBoost: a scalable tree boosting system. In:
Proceedings of the 22nd ACM SIGKDD International Conference on Knowledge
Discovery and Data Mining, pp. 785--794 (2016)

\bibitem{hochreiter1997}
Hochreiter, S., Schmidhuber, J.: Long short-term memory. Neural Comput.
\textbf{9}(8), 1735--1780 (1997)

\bibitem{vaswani2017}
Vaswani, A., et al.: Attention is all you need. In: Advances in Neural
Information Processing Systems, vol. 30, pp. 5998--6008 (2017)

\bibitem{lim2021}
Lim, B., et al.: Temporal fusion transformers
for interpretable multi-horizon time series forecasting. Int. J. Forecast.
\textbf{37}(4), 1748--1764 (2021)

\bibitem{politis1992}
Politis, D.N., Romano, J.P.: A circular block-resampling procedure for
stationary data. In: Exploring the Limits of Bootstrap, pp. 263--270. Wiley, New
York (1992)

\bibitem{engle1982}
Engle, R.F.: Autoregressive conditional heteroscedasticity with estimates of the
variance of United Kingdom inflation. Econometrica \textbf{50}(4), 987--1007
(1982)

\bibitem{carr2009}
Carr, P., Wu, L.: Variance risk premiums. Rev. Financ. Stud. \textbf{22}(3),
1311--1341 (2009)
\end{thebibliography}

\end{document}
